\chapter{Funciones usuales}\label{ch:b1:functions}

El análisis vale lo que valga el repertorio de funciones que uno
domine. A la colección heredada de la enseñanza anterior --- potencias,
exponencial, logaritmo, funciones trigonométricas --- este capítulo
añade sus funciones inversas ($\arcsin$, $\arccos$, $\arctan$) y la
familia hiperbólica. Las derivadas se usan con libertad al nivel del
volumen anterior; la teoría que hay detrás de las funciones inversas se
completa en los \cref{ch:b1:continuity,ch:b1:derivative}.

\section{Exponencial, logaritmo, potencias}

\begin{proposition}[Repaso y caracterización]\label{prop:b1:functions:expln}
$\exp \colon \R \to \intoo{0}{+\infty}$ y $\ln \colon
\intoo{0}{+\infty} \to \R$ son biyecciones recíprocas, estrictamente
crecientes, con
\[
\exp(x + y) = \exp x \exp y, \qquad \ln(xy) = \ln x + \ln y,
\qquad \exp' = \exp, \qquad \ln'(x) = \frac 1x .
\]
Además, $\exp$ es la \emph{única} función derivable $f \colon
\R \to \R$ con $f' = f$ y $f(0) = 1$.
\end{proposition}

\begin{proof}
El repaso es materia del volumen anterior. Para la unicidad, sea
$f' = f$, $f(0) = 1$, y póngase $g(x) = f(x)\,\eu^{-x}$. Entonces
$g' = f'\eu^{-x} - f\eu^{-x} = 0$, luego $g$ es constante igual a
$g(0) = 1$: $f = \exp$.
\end{proof}

\begin{definition}[Potencias de exponente real]\label{def:b1:functions:powers}
Para $x > 0$ y $\alpha \in \R$:
$\;x^\alpha = \eu^{\alpha \ln x}$\index{función potencia}. Para $a > 0$,
$a \neq 1$, el \emph{logaritmo en base $a$} es $\log_a x =
\frac{\ln x}{\ln a}$, la inversa de $x \mapsto a^x$.
\end{definition}

\begin{example}[Resolver ecuaciones exponenciales]\label{ex:b1:functions:expequations}
Resuélvase $2^x = 5^{\,x-1}$ en $\R$. Los dos miembros son positivos,
así que se toman logaritmos --- un paso reversible:
\[
x\ln2 = (x - 1)\ln5
\iff x(\ln2 - \ln5) = -\ln5
\iff x = \frac{\ln5}{\ln5 - \ln2} = \frac{\ln 5}{\ln\frac52}
\approx 1.756 .
\]
A continuación, resuélvase $x^{\sqrt2} = 3$ para $x > 0$: elévese a la
potencia $\frac1{\sqrt2}$ (es decir, aplíquese la \omterm{def:b1:logic:inj}{biyección} recíproca):
$x = 3^{1/\sqrt2} = \eu^{(\ln 3)/\sqrt2} \approx 2.175$. La idea clave:
toda ecuación que mezcle potencias se desenreda con $\ln$ y $\exp$,
porque la definición $x^\alpha = \eu^{\alpha\ln x}$ reduce toda
manipulación de potencias a la aritmética de los exponentes --- pero
solo en el dominio $x > 0$ donde esa definición vive.
\end{example}

\begin{example}[Tiempos de duplicación]\label{ex:b1:functions:doubling}
Una magnitud crece un $3\,\%$ en cada paso: tras $n$ pasos queda
multiplicada por $(1.03)^n$. ¿Cuándo se duplica? Resuélvase
$(1.03)^n \geq 2$:
\[
n\ln(1.03) \geq \ln2
\iff n \geq \frac{\ln 2}{\ln 1.03}
= \frac{0.6931}{0.02956} \approx 23.4 ,
\]
de modo que la primera duplicación ocurre en el paso $24$. (La «regla
del $72$» de los financieros, que estima el tiempo de duplicación como
$72$ dividido por la tasa en tanto por ciento, es este mismo cálculo
con la aproximación $\ln(1 + x) \approx x$, cuantificada en
el~\cref{ch:b1:taylor}.) Los procesos exponenciales se razonan mejor a
través de sus logaritmos: en esa escala el crecimiento es lineal y las
preguntas se vuelven divisiones.
\end{example}

\begin{example}[¿Cuánto ocupa $2^{2026}$?]\label{ex:b1:functions:digits}
El número de cifras decimales de un entero $N \geq 1$ es
$\floor{\log_{10} N} + 1$ (en efecto, $N$ tiene $d$ cifras exactamente
cuando $10^{d-1} \leq N < 10^d$, es decir, $d - 1 \leq \log_{10}N < d$).
Para $N = 2^{2026}$:
\[
\log_{10} 2^{2026} = 2026\,\log_{10}2
= 2026 \times 0.301030 = 609.887 ,
\]
luego $2^{2026}$ tiene $610$ cifras. La parte fraccionaria trae una
prima: $10^{0.887} \approx 7.7$, así que el número \emph{empieza} por
$7$. Una multiplicación ha respondido a una pregunta sobre un número
que nadie escribirá nunca --- los logaritmos comprimen el tamaño
multiplicativo en tamaño aditivo, que es su razón de ser histórica
(el~\cref{ex:b1:structures:expmorphism} hace precisa esa frase).
\end{example}

\begin{proposition}[Reglas de las potencias y comparación de crecimientos]\label{prop:b1:functions:powerrules}
Para $x, y > 0$ y $\alpha, \beta \in \R$:
\[
x^{\alpha+\beta} = x^\alpha x^\beta, \quad
(x^\alpha)^\beta = x^{\alpha\beta}, \quad
(xy)^\alpha = x^\alpha y^\alpha, \quad
(x^\alpha)' = \alpha\, x^{\alpha - 1}.
\]
Escala de crecimiento\index{comparación de crecimientos} cuando
$x \to +\infty$, para todos $\alpha > 0$ y $\beta > 0$:
\[
\frac{(\ln x)^\beta}{x^\alpha} \longrightarrow 0,
\qquad
\frac{x^\alpha}{\eu^{\beta x}} \longrightarrow 0 .
\]
\end{proposition}

\begin{proof}
Las identidades transcriben las de $\exp$ y $\ln$ a través de la
definición. Con detalle, la segunda (la menos evidente): $x^\alpha
> 0$ y $\ln(x^\alpha) = \alpha\ln x$ (aplíquese $\ln$ a la
definición), luego
\[
(x^\alpha)^\beta = \eu^{\beta\ln(x^\alpha)}
= \eu^{\beta\alpha\ln x} = x^{\alpha\beta} ;
\]
la primera y la tercera son los mismos desarrollos de dos líneas
mediante $\exp(u + v) = \exp u\exp v$ y $\ln(xy) = \ln x + \ln y$. La
derivada es la regla de la cadena: $(\eu^{\alpha\ln x})' =
\frac{\alpha}{x} \eu^{\alpha\ln x}$.

Comparaciones: de $\frac{\ln t}{t} \to 0$ (demostrado en el volumen
anterior), sustitúyase $t = x^{\alpha/\beta}$:
$\frac{\ln x}{x^{\alpha/\beta}} \to 0$, y elévese después a la potencia
$\beta$. Para la segunda, sustitúyase $x = \ln u$ en la primera.
\end{proof}

\begin{example}[Dos límites que conviene tener siempre a mano]\label{ex:b1:functions:xx}
¿Cuánto valen $\lim_{x \to 0^+} x^x$ y $\lim_{x \to +\infty}
x^{1/x}$? Las dos potencias se definen mediante la exponencial, así que
lo primero es resolver el exponente:
\[
x^x = \eu^{x\ln x}, \qquad
x\ln x = -\frac{\ln(1/x)}{1/x} \longrightarrow 0
\quad (x \to 0^+),
\qquad\text{so } x^x \longrightarrow \eu^0 = 1 ;
\]
y $x^{1/x} = \eu^{(\ln x)/x} \to \eu^0 = 1$ cuando $x \to +\infty$,
directamente por la \omterm{prop:b1:functions:powerrules}{comparación de crecimientos}. La idea clave: una
potencia indeterminada (de la forma $0^0$ o $\infty^0$) se trata
\emph{siempre} reescribiéndola $u^v = \eu^{v\ln u}$ y analizando el
producto $v \ln u$ --- nunca adivinando por separado a partir de la
base y del exponente. La función $x \mapsto \frac{\ln x}{x}$ que ha
decidido los dos límites se estudia a fondo en el problema del fin de
semana de este capítulo.
\end{example}

\begin{omfigure}
\begin{tikzpicture}
\begin{axis}[omaxis, width=9.5cm, height=6cm,
    xmin=0, xmax=5.2, ymin=-1.8, ymax=8.5,
    xtick={1,2,3,4,5}, ytick={2,4,6,8}]
  \addplot[omMeth, thick, domain=0.02:5, samples=200] {ln(x)}
    node[pos=0.97, below, font=\small] {$\ln x$};
  \addplot[omDef, thick, domain=0:5, samples=100] {sqrt(x)}
    node[pos=0.9, above, font=\small] {$\sqrt x$};
  \addplot[omProp, thick, domain=0:5] {x}
    node[pos=0.72, left, font=\small] {$x$};
  \addplot[omThm, thick, domain=0:2.14, samples=120] {exp(x)}
    node[pos=0.85, left, font=\small] {$\eu^x$};
\end{axis}
\end{tikzpicture}

{\small La escala de crecimiento de
la~\cref{prop:b1:functions:powerrules}: cerca del borde derecho,
$\ln x$ apenas ha pasado de $1.6$ mientras que $\eu^x$ ya se ha salido
del marco. Todos los cocientes (logaritmo entre potencia, potencia
entre exponencial) tienden a $0$ --- el dibujo solo lo insinúa; las
sustituciones de la demostración lo hacen exacto.}
\end{omfigure}

\begin{remark}[Errores frecuentes con potencias y logaritmos]\label{rem:b1:functions:pitfalls}
\begin{enumerate}
  \item \emph{Dominios.} $x^\alpha$ con $\alpha$ irracional exige
        $x > 0$; conviene evitar $(-8)^{1/3}$ y hablar de «la raíz
        cúbica real de $-8$», porque la regla
        $(x^\alpha)^\beta = x^{\alpha\beta}$ falla en silencio sobre
        los negativos: $\bigl((-8)^2\bigr)^{1/6} = 2 \neq -2$.
  \item \emph{$\sqrt{x^2} = \abs x$}, no $x$: olvidar el valor absoluto
        es la fuente clásica de soluciones negativas perdidas.
  \item \emph{$\ln(xy) = \ln x + \ln y$ exige $x, y > 0$.} Para
        argumentos negativos, $\ln(xy)$ puede estar definido sin que lo
        esté el miembro derecho.
  \item \emph{¿Qué es lo que crece?} En $x^\alpha$ varía la
        \emph{base}; en $a^x$, el \emph{exponente}. Sus crecimientos
        son radicalmente distintos
        (\cref{prop:b1:functions:powerrules}), y las expresiones
        híbridas como $x^{\ln x}$ o $x^{1/x}$ hay que reescribirlas
        como $\eu^{(\cdot)}$ antes de razonar --- como en
        el~\cref{ex:b1:functions:xx}.
\end{enumerate}
\end{remark}

\section{Funciones trigonométricas inversas}

\begin{definition}[$\arcsin$, $\arccos$, $\arctan$]\label{def:b1:functions:arc}
Las restricciones
\[
\sin \colon \intcc{-\tfrac\pi2}{\tfrac\pi2} \to \intcc{-1}{1},
\qquad
\cos \colon \intcc{0}{\pi} \to \intcc{-1}{1},
\qquad
\tan \colon \intoo{-\tfrac\pi2}{\tfrac\pi2} \to \R
\]
son biyecciones estrictamente monótonas. Sus aplicaciones inversas se
escriben $\arcsin$\index{arcoseno}, $\arccos$\index{arcocoseno} y
$\arctan$\index{arcotangente}. Así, por ejemplo, $y = \arcsin x$
($x \in \intcc{-1}{1}$) es \emph{el} ángulo de
$\intcc{-\frac\pi2}{\frac\pi2}$ cuyo seno vale $x$.
\end{definition}

\begin{omfigure}
\begin{tikzpicture}
\begin{axis}[omaxis, width=5.6cm, height=5.6cm,
    xmin=-1.25, xmax=1.25, ymin=-1.8, ymax=1.8,
    xtick={-1,1}, ytick={-1.5708,1.5708},
    yticklabels={$-\frac{\pi}{2}$,$\frac{\pi}{2}$}]
  \addplot[omDef, very thick, domain=-1:1, samples=200]
    {rad(asin(x))};
\end{axis}
\begin{scope}[xshift=5.3cm]
\begin{axis}[omaxis, width=5.6cm, height=5.6cm,
    xmin=-1.25, xmax=1.25, ymin=-0.4, ymax=3.4,
    xtick={-1,1}, ytick={1.5708,3.1416},
    yticklabels={$\frac{\pi}{2}$,$\pi$}]
  \addplot[omThm, very thick, domain=-1:1, samples=200]
    {rad(acos(x))};
\end{axis}
\end{scope}
\begin{scope}[xshift=10.6cm]
\begin{axis}[omaxis, width=5.6cm, height=5.6cm,
    xmin=-4, xmax=4, ymin=-1.9, ymax=1.9,
    xtick={-2,2}, ytick={-1.5708,1.5708},
    yticklabels={$-\frac{\pi}{2}$,$\frac{\pi}{2}$}]
  \addplot[omMeth, very thick, domain=-4:4, samples=200]
    {rad(atan(x))};
  \addplot[dashed, gray, domain=-4:4] {1.5708};
  \addplot[dashed, gray, domain=-4:4] {-1.5708};
\end{axis}
\end{scope}
\end{tikzpicture}

{\small De izquierda a derecha: $\arcsin$, $\arccos$, $\arctan$. Cada
una hereda su gráfica de la función directa restringida por reflexión
en la recta $y = x$.}
\end{omfigure}

\begin{proposition}[Derivadas]\label{prop:b1:functions:arcderiv}
En el interior de sus dominios:
\[
\arcsin' x = \frac{1}{\sqrt{1 - x^2}},
\qquad
\arccos' x = \frac{-1}{\sqrt{1 - x^2}},
\qquad
\arctan' x = \frac{1}{1 + x^2}.
\]
\end{proposition}

\begin{proof}
Anticipando la regla de derivación de la función inversa, demostrada en
el~\cref{ch:b1:derivative}: si $f$ es una \omterm{def:b1:logic:inj}{biyección} derivable con
$f' \neq 0$, entonces $(f^{-1})'(x) = \frac{1}{f'(f^{-1}(x))}$. Para
$\arcsin$: con $y = \arcsin x$,
\[
\arcsin' x = \frac{1}{\cos y}
= \frac{1}{\sqrt{1 - \sin^2 y}} = \frac{1}{\sqrt{1 - x^2}},
\]
donde $\cos y = +\sqrt{1 - \sin^2 y}$ porque $y \in
\intoo{-\frac\pi2}{\frac\pi2}$ obliga a $\cos y > 0$. Análogamente,
$\arccos' x = \frac{1}{-\sin y}$ con $\sin y > 0$ en $\intoo{0}{\pi}$,
y $\arctan' x = \frac{1}{1 + \tan^2 y} = \frac{1}{1 + x^2}$ usando
$\tan' = 1 + \tan^2$.
\end{proof}

\begin{example}[Reconocer una constante escondida]\label{ex:b1:functions:hiddenconst}
Estúdiese $g(x) = \arctan\dfrac{1 - x}{1 + x}$ en
$\intoo{-1}{+\infty}$. Regla de la cadena y un cálculo breve:
\[
g'(x) = \frac{1}{1 + \bigl(\frac{1-x}{1+x}\bigr)^2}\cdot
\frac{-(1+x) - (1-x)}{(1+x)^2}
= \frac{-2}{(1+x)^2 + (1-x)^2}
= \frac{-1}{1 + x^2} ,
\]
puesto que $(1+x)^2 + (1-x)^2 = 2 + 2x^2$. Así pues, $g' = -\arctan'$:
la función $g + \arctan$ tiene derivada nula en el \emph{intervalo}
$\intoo{-1}{+\infty}$ y, por tanto, es constante ahí; su valor en
$x = 0$ es $\arctan 1 + \arctan 0 = \frac\pi4$. Conclusión:
\[
\arctan\frac{1 - x}{1 + x} = \frac\pi4 - \arctan x
\qquad (x > -1) .
\]
En $\intoo{-\infty}{-1}$ vale el mismo cálculo de la derivada, pero la
constante es otra ($-\frac{3\pi}4$: evalúese el límite cuando
$x \to -\infty$). La idea clave: «derivada nula implica constante» es
un \omterm{def:b1:logic:statement}{enunciado} intervalo por intervalo --- exactamente la sutileza que
explotan los \cref{exo:b1:functions:6} y \cref{exo:b1:functions:10}.
\end{example}

\begin{proposition}[Identidades usuales]\label{prop:b1:functions:arcidentities}
\begin{enumerate}
  \item Para $x \in \intcc{-1}{1}$: $\arcsin x + \arccos x =
        \dfrac{\pi}{2}$.
  \item Para $x > 0$: $\arctan x + \arctan\dfrac 1x = \dfrac{\pi}{2}$
        (y $-\frac\pi2$ para $x < 0$).
  \item $\sin(\arccos x) = \cos(\arcsin x) = \sqrt{1 - x^2}$;
        $\;\tan(\arcsin x) = \dfrac{x}{\sqrt{1 - x^2}}$ para
        $\abs x < 1$.
\end{enumerate}
\end{proposition}

\begin{proof}
(1) La derivada de $x \mapsto \arcsin x + \arccos x$ es nula en
$\intoo{-1}{1}$ (\cref{prop:b1:functions:arcderiv}), luego la función
es constante ahí, igual a su valor $\frac\pi2$ en $0$; los valores en
los extremos $\pm 1$ se comprueban directamente ($\frac\pi2 + 0$ y
$-\frac\pi2 + \pi$).

(2) El mismo método en $\intoo{0}{+\infty}$: la derivada vale
$\frac{1}{1+x^2} + \frac{-1/x^2}{1 + 1/x^2} = \frac{1}{1+x^2} -
\frac{1}{x^2 + 1} = 0$, y en $x = 1$ la suma es $2\arctan 1 =
\frac\pi2$. Para $x < 0$, úsese que $\arctan$ es impar.

(3) Con $y = \arccos x \in \intcc{0}{\pi}$: $\sin y \geq 0$, luego
$\sin y = \sqrt{1 - \cos^2 y} = \sqrt{1 - x^2}$; análogamente para
$\cos(\arcsin x)$, y la fórmula de la tangente es el cociente.
\end{proof}

\begin{remark}
$\arcsin(\sin\theta) = \theta$ se cumple \emph{solo} para $\theta \in
\intcc{-\frac\pi2}{\frac\pi2}$: por ejemplo,
$\arcsin(\sin\pi) = 0 \neq \pi$. La composición en el otro sentido,
$\sin(\arcsin x) = x$, vale en todo $\intcc{-1}{1}$.
\end{remark}

\begin{example}[Volver al intervalo principal]
\label{ex:b1:functions:wrap}
Calcúlense
\[
\arctan\Bigl(\tan\frac{3\pi}4\Bigr)
\qquad\text{y}\qquad
\arctan\Bigl(\tan\frac{17\pi}5\Bigr).
\]
La receta: sustitúyase el ángulo por el único ángulo de
$\intoo{-\frac\pi2}{\frac\pi2}$ con la misma tangente, es decir,
réstese el múltiplo adecuado de $\pi$ (el período de $\tan$). Primero:
$\frac{3\pi}4 - \pi = -\frac\pi4$, luego la respuesta es $-\frac\pi4$.
Segundo: $\frac{17\pi}5 - 3\pi = \frac{2\pi}5 \in
\intoo{-\frac\pi2}{\frac\pi2}$, luego la respuesta es $\frac{2\pi}5$.
El cálculo es una división euclídea del ángulo por $\pi$ disfrazada
--- y las recetas análogas para $\arcsin$ (llevar por reflexión a
$\intcc{-\frac\pi2}{\frac\pi2}$, período $2\pi$) y $\arccos$ (llevar a
$\intcc0\pi$) mueven el~\cref{exo:b1:functions:1} y la respuesta a
trozos del~\cref{exo:b1:functions:10}.
\end{example}

\begin{example}[Sumar arcotangentes sin riesgos]\label{ex:b1:functions:arctansum}
Demostremos que
\[
\arctan\frac12 + \arctan\frac15 + \arctan\frac18 = \frac\pi4 .
\]
Dos ingredientes: la fórmula de la tangente de una suma y --- el paso
que olvidan los principiantes --- una \emph{localización} de la suma.
Primero, $\tan(u + v) = \frac{\tan u + \tan v}{1 - \tan u\tan v}$ con
$u = \arctan\frac12$, $v = \arctan\frac15$ da
\[
\tan(u + v) = \frac{\frac12 + \frac15}{1 - \frac1{10}}
= \frac{7/10}{9/10} = \frac79,
\qquad\text{y después}\qquad
\tan\Bigl(u + v + \arctan\frac18\Bigr)
= \frac{\frac79 + \frac18}{1 - \frac7{72}}
= \frac{65/72}{65/72} = 1 .
\]
Segundo, cada uno de los tres ángulos está en $\intoo0{\frac\pi4}$
(sus argumentos son menores que $1$), luego la suma está en
$\intoo0{\frac{3\pi}4}$; el único ángulo de ahí con tangente $1$ es
$\frac\pi4$. Sin la localización, la conclusión sería «$\frac\pi4$
salvo múltiplos de $\pi$» --- media demostración. La misma disciplina
en dos pasos mueve los \cref{exo:b1:functions:8}
y \cref{exo:b1:functions:12}.
\end{example}

\begin{example}[Una ecuación con arcotangentes]\label{ex:b1:functions:arctaneq}
Resuélvase $\arctan x + \arctan 2x = \dfrac\pi4$. Localícese primero:
el miembro izquierdo tiene el signo de $x$ (los dos términos lo
tienen), así que toda solución cumple $x > 0$, y entonces la suma está
en $\intoo0\pi$. Tómense tangentes (\omterm{def:b1:logic:inj}{inyectiva} en ningún intervalo de
longitud $\pi$, pero junto con la localización bastará): la fórmula de
la suma da
\[
\tan\bigl(\arctan x + \arctan 2x\bigr)
= \frac{3x}{1 - 2x^2} = 1
\iff 2x^2 + 3x - 1 = 0
\iff x = \frac{-3 \pm \sqrt{17}}4 .
\]
La raíz negativa queda descartada por la localización; para el
candidato positivo $x_0 = \frac{\sqrt{17} - 3}4 \approx 0.28$, la suma
está en $\intoo0\pi$ con tangente $1$, y el único ángulo así es
$\frac\pi4$ (en $\intoo{\frac\pi2}\pi$ la tangente es negativa): $x_0$
es la única solución. Obsérvese la forma del argumento: tomar tangentes
puede crear soluciones, nunca perderlas, así que se resuelve la
ecuación polinómica y \emph{después} se filtra por localización --- la
misma disciplina de comprobación posterior que al elevar al cuadrado
una ecuación.
\end{example}

\section{Funciones hiperbólicas}

\begin{definition}[$\cosh$, $\sinh$, $\tanh$]\label{def:b1:functions:hyperbolic}
Para $x \in \R$:
\[
\cosh x = \frac{\eu^x + \eu^{-x}}{2},
\qquad
\sinh x = \frac{\eu^x - \eu^{-x}}{2},
\qquad
\tanh x = \frac{\sinh x}{\cosh x}
\]
(\emph{coseno, seno y tangente hiperbólicos}\index{funciones
hiperbólicas}). $\cosh$ es par; $\sinh$ y $\tanh$ son impares.
\end{definition}

\begin{proposition}[Propiedades básicas]\label{prop:b1:functions:hyprules}
Para todos $x, y \in \R$:
\begin{enumerate}
  \item $\cosh^2 x - \sinh^2 x = 1$;
  \item $\cosh' = \sinh$, $\;\sinh' = \cosh$, $\;\tanh' = 1 - \tanh^2
        = \dfrac{1}{\cosh^2}$;
  \item $\cosh(x+y) = \cosh x\cosh y + \sinh x \sinh y$ y
        $\sinh(x+y) = \sinh x \cosh y + \cosh x \sinh y$;
  \item $\sinh$ es una \omterm{def:b1:logic:inj}{biyección} estrictamente creciente de $\R$ sobre
        $\R$; $\cosh$ restringido a $\R_+$ es una \omterm{def:b1:logic:inj}{biyección}
        estrictamente creciente sobre $\intco{1}{+\infty}$; $\tanh$ es
        una \omterm{def:b1:logic:inj}{biyección} estrictamente creciente de $\R$ sobre $\intoo{-1}{1}$.
\end{enumerate}
\end{proposition}

\begin{proof}
(1) $\bigl(\frac{\eu^x + \eu^{-x}}{2}\bigr)^2 - \bigl(\frac{\eu^x -
\eu^{-x}}{2}\bigr)^2 = \frac{(\eu^{2x} + 2 + \eu^{-2x}) - (\eu^{2x} -
2 + \eu^{-2x})}{4} = 1$.

(2) Se derivan las fórmulas de definición; para $\tanh$, la regla del
cociente da $\frac{\cosh^2 - \sinh^2}{\cosh^2}$, que es a la vez
$1 - \tanh^2$ y $\frac{1}{\cosh^2}$ por (1).

(3) Se desarrollan los miembros derechos con las definiciones. Con
detalle para la primera fórmula:
\[
\cosh x\cosh y + \sinh x\sinh y
= \frac{(\eu^x + \eu^{-x})(\eu^y + \eu^{-y})
+ (\eu^x - \eu^{-x})(\eu^y - \eu^{-y})}{4} ;
\]
de los ocho productos, los cuatro «mixtos» ($\eu^{x-y}$,
$\eu^{y-x}$) se cancelan por parejas, mientras que $\eu^{x+y}$ y
$\eu^{-(x+y)}$ aparecen cada uno dos veces: el total es
$\frac{2\eu^{x+y} + 2\eu^{-(x+y)}}{4} = \cosh(x+y)$. La fórmula del
seno es idéntica, sobreviviendo los términos mixtos en lugar de los
otros.

(4) $\sinh' = \cosh \geq 1 > 0$, luego $\sinh$ es estrictamente
creciente, con límites $\pm\infty$ (término dominante
$\pm\frac{\eu^{\abs x}}{2}$); el \omterm{def:b1:logic:statement}{enunciado} sobre la \omterm{def:b1:logic:inj}{biyección} se sigue
entonces del teorema del valor intermedio (usado al nivel del volumen
anterior; tratamiento sistemático en el~\cref{ch:b1:continuity}). En
$\R_+$, $\cosh' = \sinh \geq 0$ se anula solo en $0$: estrictamente
creciente desde $\cosh 0 = 1$ hasta $+\infty$. Y $\tanh' > 0$ con
límites $\pm 1$ en $\pm\infty$: con detalle,
detail,
\[
\tanh x = \frac{\eu^x - \eu^{-x}}{\eu^x + \eu^{-x}}
= \frac{1 - \eu^{-2x}}{1 + \eu^{-2x}}
\xrightarrow[x \to +\infty]{} 1 ,
\]
tras dividir numerador y denominador por $\eu^x$, y la imparidad da el
límite $-1$ en $-\infty$; la función estrictamente creciente $\tanh$
envía por tanto $\R$ sobre $\intoo{-1}1$.
\end{proof}

\begin{omfigure}
\begin{tikzpicture}
\begin{axis}[omaxis, width=6.4cm, height=6cm,
    xmin=-2.6, xmax=2.6, ymin=-3.6, ymax=3.9,
    xtick={-2,-1,1,2}, ytick={-3,-2,-1,1,2,3}]
  \addplot[omDef, very thick, domain=-2.05:2.05, samples=120]
    {cosh(x)} node[pos=0.1, right, font=\small] {$\cosh$};
  \addplot[omThm, very thick, domain=-2.05:2.05, samples=120]
    {sinh(x)} node[pos=0.93, above left, font=\small] {$\sinh$};
  \addplot[dashed, gray, domain=0:2.5] {exp(x)/2};
\end{axis}
\begin{scope}[xshift=7.4cm]
\begin{axis}[omaxis, width=6.4cm, height=6cm,
    xmin=-3.2, xmax=3.2, ymin=-1.5, ymax=1.5,
    xtick={-2,-1,1,2}, ytick={-1,1}]
  \addplot[omMeth, very thick, domain=-3:3, samples=150]
    {tanh(x)} node[pos=0.9, above left, font=\small] {$\tanh$};
  \addplot[dashed, gray, domain=-3:3] {1};
  \addplot[dashed, gray, domain=-3:3] {-1};
\end{axis}
\end{scope}
\end{tikzpicture}

{\small A la izquierda: $\cosh$ y $\sinh$, asintóticamente pegados a
$\frac{\eu^x}{2}$ (a trazos). A la derecha: $\tanh$, creciente de $-1$
a $1$.}
\end{omfigure}

\begin{remark}[¿Por qué «hiperbólicas»?]
El punto $(\cosh t, \sinh t)$ recorre la rama $x > 0$ de la hipérbola
$x^2 - y^2 = 1$ (por la~\cref{prop:b1:functions:hyprules}~(1)),
exactamente como $(\cos t, \sin t)$ recorre la circunferencia
$x^2 + y^2 = 1$. Toda identidad circular tiene una hermana
hiperbólica, con cambios de signo gobernados por
$\sinh^2 \leftrightarrow -\sin^2$.
\end{remark}

\begin{example}[La fórmula de adición de $\tanh$]\label{ex:b1:functions:tanhadd}
Dividiendo por $\cosh x\cosh y$ las dos fórmulas de adición de
la~\cref{prop:b1:functions:hyprules}~(3):
\[
\tanh(x + y)
= \frac{\sinh x\cosh y + \cosh x\sinh y}
{\cosh x\cosh y + \sinh x\sinh y}
= \frac{\tanh x + \tanh y}{1 + \tanh x\,\tanh y} ,
\]
la hermana hiperbólica de la fórmula de adición de la tangente --- con
un $+$ donde la trigonometría tiene un $-$. Un dividendo: como
$\abs{\tanh} < 1$, el miembro derecho es una regla de «suma de
velocidades» que nunca sale de $\intoo{-1}1$: si $u, v \in \intoo{-1}1$,
entonces $\frac{u + v}{1 + uv} \in \intoo{-1}1$ también (escríbanse
$u = \tanh a$, $v = \tanh b$, posible por \omterm{def:b1:logic:inj}{biyectividad}, y léase la
fórmula al revés). Comprobarlo por álgebra, sin \omterm{def:b1:functions:hyperbolic}{funciones
hiperbólicas}, es un ejercicio algo penoso; parametrizar con $\tanh$ lo
convierte en una línea --- la misma estrategia que las funciones
circulares brindan para la circunferencia unidad.
\end{example}

\begin{proposition}[Funciones hiperbólicas inversas]\label{prop:b1:functions:invhyp}
Las inversas que garantiza la~\cref{prop:b1:functions:hyprules}~(4)
tienen forma cerrada:
\[
\operatorname{arsinh} x = \ln\bigl(x + \sqrt{x^2 + 1}\bigr)
\ (x \in \R),
\qquad
\operatorname{arcosh} x = \ln\bigl(x + \sqrt{x^2 - 1}\bigr)
\ (x \geq 1),
\]
\[
\operatorname{artanh} x = \frac 12 \ln\frac{1 + x}{1 - x}
\ (\abs x < 1),
\]
con derivadas $\frac{1}{\sqrt{x^2+1}}$, $\frac{1}{\sqrt{x^2-1}}$
($x > 1$) y $\frac{1}{1 - x^2}$, respectivamente.
\end{proposition}

\begin{proof}
Para $\operatorname{arsinh}$: resuélvase $x = \sinh y = \frac{\eu^y -
\eu^{-y}}{2}$. Poniendo $u = \eu^y > 0$: $u^2 - 2xu - 1 = 0$, luego
$u = x + \sqrt{x^2 + 1}$ (la raíz $x - \sqrt{x^2+1}$ es negativa), y
$y = \ln(x + \sqrt{x^2+1})$. Las otras dos son cálculos idénticos
($u^2 - 2xu + 1 = 0$ para $\operatorname{arcosh}$, quedándose con la
raíz $\geq 1$; una reordenación de dos líneas para
$\operatorname{artanh}$). Derivadas: se derivan las expresiones
logarítmicas; por ejemplo,
\[
\bigl(\ln(x + \sqrt{x^2+1})\bigr)'
= \frac{1 + \frac{x}{\sqrt{x^2+1}}}{x + \sqrt{x^2+1}}
= \frac{1}{\sqrt{x^2 + 1}} . \qedhere
\]
\end{proof}

\begin{example}[Las formas cerradas en acción]\label{ex:b1:functions:arcoshtwo}
La solución de $\cosh t = 2$ con $t \geq 0$ es, por la forma cerrada,
$t = \operatorname{arcosh} 2 = \ln(2 + \sqrt3) \approx 1.317$; la otra
solución es $-t$, por paridad --- y, en efecto,
$\ln(2 - \sqrt3) = \ln\frac1{2 + \sqrt3} = -\ln(2 + \sqrt3)$: las dos
raíces $u = \eu^t$ de la cuadrática $u^2 - 4u + 1 = 0$ son inversas la
una de la otra, como exige su producto $1$ (Vieta). Este cálculo
mínimo muestra el patrón general: las ecuaciones hiperbólicas se
convierten en cuadráticas en $\eu^t$, y la simetría $t \mapsto -t$
aparece como la simetría $u \mapsto 1/u$ de la cuadrática --- conviene
recordarlo al resolver el~\cref{exo:b1:functions:7}.
\end{example}

\begin{method}[Elegir la primitiva adecuada]\label{met:b1:functions:primitive}
Los tres patrones de derivada que hay que memorizar para la integración
(\cref{ch:b1:integration}):
\[
\int \frac{\dd x}{1 + x^2} = \arctan x + C,
\qquad
\int \frac{\dd x}{\sqrt{1 - x^2}} = \arcsin x + C,
\qquad
\int \frac{\dd x}{\sqrt{x^2 + 1}} = \operatorname{arsinh} x + C,
\]
y, para una cuadrática cualquiera, se reduce a estos completando
cuadrados y reescalando.
\end{method}

\begin{remark}[Dónde se usa este capítulo]\label{rem:b1:functions:whereused}
Este capítulo es el vocabulario de trabajo de todo el análisis que
viene. La escala de crecimiento de
la~\cref{prop:b1:functions:powerrules} decide cuestiones de
convergencia en los \cref{ch:b1:seq,ch:b1:series}; las fórmulas de
derivación de las \cref{prop:b1:functions:arcderiv}
y \cref{prop:b1:functions:invhyp} son las primitivas que más falta
hacen en el~\cref{ch:b1:integration}, vía
el~\cref{met:b1:functions:primitive}; las \omterm{def:b1:functions:hyperbolic}{funciones hiperbólicas}
parametrizan las soluciones de la ecuación $y'' = y$ en
el~\cref{ch:b1:diffeq} exactamente como las circulares parametrizan las
de $y'' = -y$. La caracterización de $\exp$ por $f' = f$, $f(0) = 1$
(\cref{prop:b1:functions:expln}) es la semilla unidimensional de la
teoría de las ecuaciones diferenciales lineales, y la catenaria
$y = \cosh x$ vuelve entre las curvas planas del~\cref{ch:b1:curves}.
\end{remark}

\begin{remark}[Interludio: el programa de la función inversa]\label{rem:b1:functions:interlude}
Este capítulo ha ejecutado cuatro veces el mismo programa: restringir
una función hasta que sea una \omterm{def:b1:logic:inj}{biyección} estrictamente monótona, dar
nombre a la inversa y transportar por ella todas las fórmulas. Lo que
se ha \emph{usado} en cada paso --- que una función continua y
estrictamente monótona en un intervalo es una \omterm{def:b1:logic:inj}{biyección} sobre un
intervalo, y que su inversa es continua y después derivable fuera de
los puntos críticos --- se ha tomado a crédito del volumen anterior. La
deuda se salda dentro de este volumen: el~\cref{ch:b1:continuity}
demuestra el \omterm{def:b1:logic:statement}{enunciado} sobre la \omterm{def:b1:logic:inj}{biyección} (el teorema de la inversa
monótona, apoyado en el teorema del valor intermedio), y
el~\cref{ch:b1:derivative} demuestra la regla de derivación
$(f^{-1})' = 1/(f' \circ f^{-1})$ que ha producido en silencio todas
las fórmulas de las \cref{prop:b1:functions:arcderiv}
y \cref{prop:b1:functions:invhyp}. Leer esos capítulos teniendo en
mente $\arcsin$ y $\operatorname{arsinh}$ --- como los ejemplos
resueltos que la teoría se construyó para justificar --- es la manera
prevista de sortear la aparente circularidad.
\end{remark}

\section{Ejercicios}

\begin{exercise}[$\star$]\label{exo:b1:functions:1}
Calcúlense sin calculadora: $\arcsin\frac{\sqrt 3}{2}$;
$\;\arccos\bigl(-\frac 12\bigr)$; $\;\arctan(-1)$;
$\;\arcsin\bigl(\sin\frac{5\pi}{6}\bigr)$;
$\;\arccos\bigl(\cos\frac{7\pi}{4}\bigr)$.
\end{exercise}

\begin{exercise}[$\star$]\label{exo:b1:functions:2}
Dese el dominio de definición de $f(x) = \arcsin(2x - 1)$ y calcúlese
$f'$ donde exista. Las mismas preguntas para $g(x) = \arctan\sqrt{x}$.
\end{exercise}

\begin{exercise}[$\star$]\label{exo:b1:functions:3}
Demuéstrese que, para todo $x \in \R$, $\cosh x + \sinh x = \eu^x$ y
$(\cosh x + \sinh x)^n = \cosh nx + \sinh nx$ para todo $n \in \Z$
(una fórmula de De Moivre hiperbólica).
\end{exercise}

\begin{exercise}[$\star$]\label{exo:b1:functions:4}
Simplifíquense $\cos(2\arcsin x)$ y $\sin(2\arctan x)$ hasta obtener
expresiones algebraicas de $x$.
\end{exercise}

\begin{exercise}[$\star$]\label{exo:b1:functions:5}
Ordénense, para $x$ grande, las funciones
$x^{100}$, $\;\eu^{\sqrt x}$, $\;(\ln x)^{1000}$, $\;\eu^{x/100}$,
$\;x^{\ln x}$, del crecimiento más lento al más rápido, con
justificaciones basadas en la~\cref{prop:b1:functions:powerrules}.
\end{exercise}

\begin{exercise}[$\star\star$]\label{exo:b1:functions:6}
Estúdiese la función $f(x) = \arctan\dfrac{2x}{1 - x^2}$ en su dominio:
calcúlese $f'$, compárese con $(2\arctan x)'$ y exprésese $f(x)$ en
términos de $\arctan x$ en cada uno de los tres intervalos del dominio.
\end{exercise}

\begin{exercise}[$\star\star$]\label{exo:b1:functions:7}
Resuélvanse en $\R$: $\;\cosh x = 2$; después $5\cosh x - 4 \sinh x = 3$
(exprésense las soluciones con logaritmos). \emph{Indicación para la
segunda: escríbase todo con $u = \eu^x$.}
\end{exercise}

\begin{exercise}[$\star\star$]\label{exo:b1:functions:8}
Demuéstrese la identidad $\arctan\frac{1}{2} + \arctan\frac{1}{3} =
\frac{\pi}{4}$ y después la fórmula de Machin\index{fórmula de Machin}
\[
4\arctan\frac 15 - \arctan\frac{1}{239} = \frac{\pi}{4}.
\]
\emph{Indicación: calcúlese la tangente de los dos miembros con la
fórmula de adición y contrólese en qué intervalo cae cada uno.}
\end{exercise}

\begin{exercise}[$\star\star$]\label{exo:b1:functions:9}
Demuéstrese que, para todo $x \geq 0$, $\;x - \dfrac{x^3}{6} \leq \sin x
\leq x$ \emph{(estúdiense las derivadas sucesivas de las diferencias)},
y dedúzcase que $\lim_{x \to 0^+} \frac{x - \sin x}{x^3} = \frac 16$ es
plausible --- el límite en sí se establece en
el~\cref{ch:b1:taylor}.
\end{exercise}

\begin{exercise}[$\star\star\star$]\label{exo:b1:functions:10}
Para $x \in \intcc{-1}{1}$, póngase
$h(x) = \arcsin\bigl(2x\sqrt{1 - x^2}\,\bigr) - 2\arcsin x$. Cabría
esperar $h = 0$ a partir de la identidad $\sin 2y = 2\sin y\cos y$
--- pero $h$ \emph{no} es idénticamente nula. Calcúlese $h'$ en los
intervalos abiertos donde exista, evalúese $h$ en puntos bien elegidos
y dese la descripción completa de $h$ como función constante a trozos
en $\intcc{-1}{1}$.
\end{exercise}

\begin{exercise}[$\star\star\star$]\label{exo:b1:functions:11}
(Gudermanniana) Sea $g(x) = \arctan(\sinh x)$ para $x \in \R$.
Demuéstrese que $g$ es una \omterm{def:b1:logic:inj}{biyección} impar y estrictamente creciente de
$\R$ sobre $\intoo{-\frac\pi2}{\frac\pi2}$, que $g'(x) =
\frac{1}{\cosh x}$ y que $\tan g(x) = \sinh x$, $\;\sin g(x) = \tanh x$,
$\;\cos g(x) = \frac{1}{\cosh x}$: la función $g$ enlaza la
trigonometría circular con la hiperbólica sin números complejos.
\end{exercise}

\begin{exercise}[$\star\star$]\label{exo:b1:functions:12}
Demuéstrese que
\[
\arctan 1 + \arctan 2 + \arctan 3 = \pi .
\]
\emph{Indicación: calcúlese primero $\arctan 2 + \arctan 3$,
localizando la suma como en el~\cref{ex:b1:functions:arctansum};
cuidado, aquí el denominador de la fórmula de adición es negativo.}
\end{exercise}

\section{Problema: la ecuación $x^y = y^x$}

\begin{problem}\label{pb:b1:functions:1}
¿Qué pares de números positivos cumplen $x^y = y^x$? Todo el mundo
conoce un ejemplo de apariencia accidental, $2^4 = 4^2 = 16$; este
problema muestra que en él no hay nada de accidental. Toda la ecuación
está gobernada por las variaciones de una única función
\[
f(t) = \frac{\ln t}{t} \qquad (t > 0),
\]
cuyo estudio proporciona: el \omterm{def:b1:logic:sets}{conjunto} completo de soluciones (una
diagonal más una rama curva que pasa por $(\eu, \eu)$), una
parametrización racional de la rama, el hecho de que $(2, 4)$ es su
\emph{único} punto entero y la clasificación de todos sus puntos
racionales --- además de, como dividendos, la comparación de $\eu^\pi$
con $\pi^\eu$ y la convergencia monótona de
$\bigl(1 + \frac1n\bigr)^n$ hacia $\eu$.
\index{función potencia}

\medskip
\noindent\textbf{Parte I --- La función $f(t) = \ln t / t$.}
\begin{enumerate}
  \item Justifíquese que $f$ es derivable en $\intoo0{+\infty}$,
        calcúlese $f'$ y hágase la tabla de variación: $f$ crece
        estrictamente en $\intoc0\eu$ y decrece estrictamente en
        $\intco\eu{+\infty}$, con máximo $f(\eu) = \frac1\eu$.
  \item Determínense los límites de $f$ en $0^+$ y en $+\infty$
        (\cref{prop:b1:functions:powerrules}), el signo de $f$
        (negativo en $\intoo01$, nulo en $1$, positivo más allá) y
        esbócese la gráfica.
  \item Compruébese por cálculo directo que $f(2) = f(4)$. (Téngase
        presente esta igualdad: de ella nace todo el problema.)
  \item Sea $c \in \R$. Discútase, según el valor de $c$, el número de
        soluciones de $f(t) = c$: exactamente una para $c \leq 0$;
        exactamente dos (una en $\intoo1\eu$ y otra en
        $\intoo\eu{+\infty}$) para $0 < c < \frac1\eu$; exactamente una
        para $c = \frac1\eu$; ninguna para $c > \frac1\eu$.
  \item Dedúzcase que $\eu^x > x^\eu$ para todo $x > 0$ con
        $x \neq \eu$ y, en particular, decídase cuál es mayor, $\eu^\pi$
        o $\pi^\eu$.
\end{enumerate}

\medskip
\noindent\textbf{Parte II --- La ecuación y su curva.} En esta parte,
$x, y > 0$.
\begin{enumerate}[resume]
  \item Pruébese que $x^y = y^x \iff f(x) = f(y)$.
  \item Dedúzcase la estructura del \omterm{def:b1:logic:sets}{conjunto} de soluciones: todos los
        pares diagonales $(x, x)$; y los pares \emph{no triviales}
        ($x \neq y$), que cumplen que las dos coordenadas son $> 1$ y
        que una está en $\intoo1\eu$ mientras la otra está en
        $\intoo\eu{+\infty}$.
  \item Parametrícense los pares no triviales: escribiendo $y = tx$ con
        $t > 0$, $t \neq 1$, pruébese que $x^y = y^x$ obliga a
        \[
        x(t) = t^{\frac1{t-1}},
        \qquad
        y(t) = t\,x(t) = t^{\frac{t}{t-1}},
        \]
        y que, recíprocamente, todo par así es solución.
  \item Compruébese que $t = 2$ da $(2, 4)$ y demuéstrese la simetría
        $x(1/t) = y(t)$, $y(1/t) = x(t)$: invertir el parámetro
        intercambia las dos coordenadas.
  \item Determínense los límites de $x(t)$ e $y(t)$ cuando $t \to 1$
        (los dos tienden a $\eu$), cuando $t \to +\infty$ ($x \to 1$,
        $y \to +\infty$) y cuando $t \to 0^+$ ($x \to +\infty$,
        $y \to 1$). Descríbase la rama resultante: una curva asintótica
        a las rectas $x = 1$ e $y = 1$, que corta a la diagonal en
        $(\eu, \eu)$.
  \item Pruébese que $t \mapsto x(t)$ es estrictamente decreciente en
        $\intoo1{+\infty}$. \emph{(Estúdiese $h(t) = 1 - \frac1t - \ln
        t$, cuyo signo controla la derivada de $\ln x(t) =
        \frac{\ln t}{t - 1}$.)}
  \item Conclúyase que las soluciones no triviales definen una
        \omterm{def:b1:logic:inj}{biyección} estrictamente decreciente $\varphi \colon \intoo1\eu
        \to \intoo\eu{+\infty}$ con $x^{\varphi(x)} = \varphi(x)^x$, y
        que $\varphi$ es una involución de la rama:
        $\varphi(\varphi(x)) = x$ donde ambos miembros estén definidos.
\end{enumerate}

\medskip
\noindent\textbf{Parte III --- Puntos enteros y racionales.}
\begin{enumerate}[resume]
  \item Demuéstrese que $(2, 4)$ y $(4, 2)$ son las únicas soluciones
        \emph{enteras} no triviales de $x^y = y^x$.
  \item Para $n \in \N^*$, aplíquese la parametrización con
        $t = 1 + \frac1n$ y pruébese que
        \[
        x_n = \Bigl(1 + \frac1n\Bigr)^{n},
        \qquad
        y_n = \Bigl(1 + \frac1n\Bigr)^{n+1}
        \]
        es una solución \emph{racional} no trivial para todo $n$, con
        $(x_1, y_1) = (2, 4)$.
  \item Recíprocamente, sea $(x, y)$ una solución racional no trivial
        con $y > x$, y escríbase $t = y/x = 1 + \frac rs$ como fracción
        irreducible ($r, s \in \N^*$). Usando $x = t^{s/r}$ y
        admitiendo la unicidad de la factorización en primos (conocida
        de la enseñanza anterior; demostrada en el~\cref{ch:b1:arith}),
        pruébese que ser $x$ racional obliga a que $s$ y $s + r$ sean
        ambos potencias $r$-ésimas de enteros.
  \item Pruébese con el teorema del binomio que $a^r - b^r = r$ no
        tiene soluciones enteras $a > b \geq 1$ cuando $r \geq 2$, y
        conclúyase: las soluciones racionales de $x^y = y^x$ son
        exactamente los pares $(x_n, y_n)$ de la pregunta 14 (y los
        pares intercambiados).
  \item Verifíquese numéricamente el caso $n = 2$: calcúlense $f(9/4)$
        y $f(27/8)$ con cuatro decimales y compruébese que coinciden.
\end{enumerate}

\medskip
\noindent\textbf{Parte IV --- Dividendos.}
\begin{enumerate}[resume]
  \item Úsense las preguntas 7 y 11 para demostrar, sin ningún cálculo
        adicional, que la sucesión $x_n = \bigl(1 + \frac1n\bigr)^n$ es
        estrictamente creciente con $x_n < \eu$, que
        $y_n = \bigl(1 + \frac1n\bigr)^{n+1}$ es estrictamente
        decreciente con $y_n > \eu$, y que ambas convergen a $\eu$. (El
        parámetro $t_n = 1 + \frac1n$ decrece hacia $1$.)
  \item Establézcase la regla general de comparación para $1 < a < b$:
        si $\eu \leq a$, entonces $a^b > b^a$; si $b \leq \eu$,
        entonces $a^b < b^a$; y véase con los dos ejemplos $(2, 3)$ y
        $(2, 5)$ que en el caso mixto $a < \eu < b$ se dan realmente
        los dos resultados.
  \item Supóngase que la involución $\varphi$ de la pregunta 12 es
        derivable en $\eu$ (lo es). Derívese la identidad
        $\varphi(\varphi(x)) = x$ en $x = \eu$ y dedúzcase
        $\varphi'(\eu) = -1$: la rama corta a la diagonal
        perpendicularmente.
  \item Hállese en forma cerrada el único par solución con $y = 3x$ y
        compruébese numéricamente con cuatro decimales mediante $f$.
  \item Entre los números $\sqrt2$, $\sqrt[3]3$, $\sqrt[4]4$,
        $\sqrt[5]5$, determínese el mayor e identifíquense los dos que
        son iguales. \emph{(Compárense $n^{1/n} = \eu^{f(n)}$.)}
\end{enumerate}

\medskip
\noindent\textbf{Parte V --- Síntesis.}
\begin{enumerate}[resume]
  \item Descríbase el \omterm{def:b1:logic:sets}{conjunto} completo de soluciones de $x^y = y^x$
        en el cuadrante $x, y > 0$ --- diagonal más rama, su
        intersección $(\eu, \eu)$, las asíntotas, el punto entero
        $(2, 4)$, los puntos racionales que se acumulan en
        $(\eu, \eu)$ --- de una forma que se pueda esbozar de memoria.
  \item ¿Dónde ha usado exactamente el problema: (i) las
        comparaciones de crecimiento de
        la~\cref{prop:b1:functions:powerrules}; (ii) el teorema del
        valor intermedio (a través de los \omterm{def:b1:logic:statement}{enunciados} sobre biyecciones);
        (iii) la unicidad admitida de la factorización en primos? Una
        frase para cada uno.
  \item Moraleja, en un párrafo breve: una sola tabla de variación ha
        resuelto una ecuación con dos incógnitas, ha clasificado sus
        puntos racionales y ha demostrado la convergencia monótona de
        $\bigl(1 + \frac1n\bigr)^n$ --- coméntese esta economía y
        nómbrese dónde se industrializa cada hilo más adelante en el
        volumen (\cref{ch:b1:seq} para la sucesión,
        \cref{ch:b1:derivative} para las tablas de variación,
        \cref{ch:b1:taylor} para la precisión que a la tabla le falta).
\end{enumerate}
\end{problem}
